\section{其他形式的曲线、曲面积分}

\begin{problem}{散度的积分表示推导}{Divergence-Integral-Cylindrical}
    利用散度的积分表示，推导出在柱坐标系下的散度。
\end{problem}

\begin{solution}
    散度的积分表示定义为：
    \begin{equation}
        \operatorname{div} \symbfit{A} = \lim_{\Delta V \to 0} \frac{1}{\Delta V} \oiint_{S} \symbfit{A} \cdot \d \symbfit{S},
    \end{equation}
    其中 $S$ 为包围体积元 $\Delta V$ 的闭合曲面。在柱坐标系 $(r, \theta, z)$ 中，设向量场表示为 $\symbfit{A} = A_r \symbfit{e}_r + A_\theta \symbfit{e}_\theta + A_z \symbfit{e}_z$。选取由坐标面构成的微小“六面体”体积元 $\Delta V$，其边界范围在径向为 $r$ 到 $r+\Delta r$、在角向为 $\theta$ 到 $\theta+\Delta \theta$ 以及在轴向为 $z$ 到 $z+\Delta z$。该体积元的大小可近似为：
    \begin{equation}
        \Delta V \approx r \Delta r \Delta \theta \Delta z.
    \end{equation}

    我们分别计算向量场 $\symbfit{A}$ 穿过该体积元六个面的向外总通量。在径向（$r$ 方向）的两个面上，法向面微元大小分别为 $r \Delta \theta \Delta z$ 和 $(r+\Delta r) \Delta \theta \Delta z$，向外的净通量为：
    \begin{equation}
        \begin{aligned}
            \Phi_r & \approx A_r(r+\Delta r, \theta, z) (r+\Delta r) \Delta \theta \Delta z - A_r(r, \theta, z) r \Delta \theta \Delta z \\
                   & \approx \frac{\partial (r A_r)}{\partial r} \Delta r \Delta \theta \Delta z.
        \end{aligned}
    \end{equation}

    在角向（$\theta$ 方向）的两个面上，法向面微元大小均为 $\Delta r \Delta z$，向外的净通量为：
    \begin{equation}
        \begin{aligned}
            \Phi_\theta & \approx A_\theta(r, \theta+\Delta \theta, z) \Delta r \Delta z - A_\theta(r, \theta, z) \Delta r \Delta z \\
                        & \approx \frac{\partial A_\theta}{\partial \theta} \Delta \theta \Delta r \Delta z.
        \end{aligned}
    \end{equation}

    在轴向（$z$ 方向）的两个面上，法向面微元大小均为 $r \Delta r \Delta \theta$，向外的净通量为：
    \begin{equation}
        \begin{aligned}
            \Phi_z & \approx A_z(r, \theta, z+\Delta z) r \Delta r \Delta \theta - A_z(r, \theta, z) r \Delta r \Delta \theta \\
                   & \approx \frac{\partial A_z}{\partial z} \Delta z \, r \Delta r \Delta \theta.
        \end{aligned}
    \end{equation}

    将上述三个方向的通量相加，即可得到穿过闭合曲面 $S$ 的向外总通量：
    \begin{equation}
        \oiint_{S} \symbfit{A} \cdot \d \symbfit{S} = \Phi_r + \Phi_\theta + \Phi_z \approx \left( \frac{\partial (r A_r)}{\partial r} + \frac{\partial A_\theta}{\partial \theta} + r \frac{\partial A_z}{\partial z} \right) \Delta r \Delta \theta \Delta z.
    \end{equation}

    将其代入散度的积分表示定义式中，并除以体积微元 $\Delta V = r \Delta r \Delta \theta \Delta z$，在取极限 $\Delta V \to 0$ 的过程中消去高阶无穷小，最终得到柱坐标系下的散度公式：
    \begin{equation}
        \boxed{\operatorname{div} \symbfit{A} = \frac{1}{r} \frac{\partial (r A_r)}{\partial r} + \frac{1}{r} \frac{\partial A_\theta}{\partial \theta} + \frac{\partial A_z}{\partial z}}
    \end{equation}
\end{solution}

\begin{problem}{梯度的积分表示推导}{Gradient-Integral-Spherical}
    利用梯度的积分表示，推导出在球坐标系下的梯度。
\end{problem}

\begin{solution}
    梯度的积分表示定义为：
    \begin{equation}
        \operatorname{grad} f = \lim_{\Delta V \to 0} \frac{1}{\Delta V} \oiint_{S} f \d \symbfit{S},
    \end{equation}
    其中 $\d \symbfit{S} = \symbfit{n} \d S$ 为外法向面积微元。在球坐标系 $(r, \theta, \phi)$ 中，选取由坐标面构成的微小体积元，边界范围分别为 $r$ 到 $r+\Delta r$、$\theta$ 到 $\theta+\Delta \theta$、$\phi$ 到 $\phi+\Delta \phi$，其体积可近似表示为：
    \begin{equation}
        \Delta V \approx r^2 \sin\theta \Delta r \Delta \theta \Delta \phi.
    \end{equation}

    我们将闭合曲面积分拆解为径向（$r$）、极角（$\theta$）和方位角（$\phi$）三对相对面上的面积分近似之和。在径向的两个面上，面积微元大小约为 $r^2 \sin\theta \Delta \theta \Delta \phi$，外法向量为 $\pm \symbfit{e}_r$，其积分贡献为：
    \begin{equation}
        \begin{aligned}
            I_r & \approx \left( f(r+\Delta r) \symbfit{e}_r(r+\Delta r)^2 - f(r) \symbfit{e}_r r^2 \right) \sin\theta \Delta \theta \Delta \phi \\
                & \approx \frac{\partial}{\partial r} (f r^2 \symbfit{e}_r) \sin\theta \Delta r \Delta \theta \Delta \phi.
        \end{aligned}
    \end{equation}

    在极角的两个面上，面积微元大小约为 $r \sin\theta \Delta r \Delta \phi$，外法向量为 $\pm \symbfit{e}_\theta$，积分贡献为：
    \begin{equation}
        I_\theta \approx \frac{\partial}{\partial \theta} (f \sin\theta \symbfit{e}_\theta) r \Delta r \Delta \theta \Delta \phi.
    \end{equation}

    在方位角的两个面上，面积微元大小约为 $r \Delta r \Delta \theta$，外法向量为 $\pm \symbfit{e}_\phi$，积分贡献为：
    \begin{equation}
        I_\phi \approx \frac{\partial}{\partial \phi} (f \symbfit{e}_\phi) r \Delta r \Delta \theta \Delta \phi.
    \end{equation}

    在球坐标系中，局部基向量对各坐标变量的偏导数满足如下几何关系：
    \begin{equation}
        \label{eq:base-vector-deriv}
        \frac{\partial \symbfit{e}_r}{\partial r} = \symbfit{0}, \quad \frac{\partial \symbfit{e}_\theta}{\partial \theta} = -\symbfit{e}_r, \quad \frac{\partial \symbfit{e}_\phi}{\partial \phi} = -\sin\theta \symbfit{e}_r - \cos\theta \symbfit{e}_\theta.
    \end{equation}

    利用\zcref{eq:base-vector-deriv}中基向量的导数公式，利用乘积法则分别展开 $I_r$、$I_\theta$、$I_\phi$ 中的偏导数：
    \begin{equation}
        \begin{aligned}
            I_r      & \approx \symbfit{e}_r \sin\theta \left( r^2 \frac{\partial f}{\partial r} + 2r f \right) \Delta r \Delta \theta \Delta \phi,                                                                                                      \\
            I_\theta & \approx r \left( \frac{\partial f}{\partial \theta} \symbfit{e}_\theta \sin\theta + f \frac{\partial \symbfit{e}_\theta}{\partial \theta} \sin\theta + f \symbfit{e}_\theta \cos\theta \right) \Delta r \Delta \theta \Delta \phi \\
                     & = r \left( \frac{\partial f}{\partial \theta} \symbfit{e}_\theta \sin\theta - f \symbfit{e}_r \sin\theta + f \symbfit{e}_\theta \cos\theta \right) \Delta r \Delta \theta \Delta \phi,                                            \\
            I_\phi   & \approx r \left( \frac{\partial f}{\partial \phi} \symbfit{e}_\phi + f \frac{\partial \symbfit{e}_\phi}{\partial \phi} \right) \Delta r \Delta \theta \Delta \phi                                                                 \\
                     & = r \left( \frac{\partial f}{\partial \phi} \symbfit{e}_\phi - f \sin\theta \symbfit{e}_r - f \cos\theta \symbfit{e}_\theta \right) \Delta r \Delta \theta \Delta \phi.
        \end{aligned}
    \end{equation}

    将上述三部分积分相加，所有包含纯代数 $f$ 本身的项恰好相互抵消：
    \begin{equation}
        \oiint_{S} f \d \symbfit{S} = I_r + I_\theta + I_\phi \approx \left( r^2 \sin\theta \frac{\partial f}{\partial r} \symbfit{e}_r + r \sin\theta \frac{\partial f}{\partial \theta} \symbfit{e}_\theta + r \frac{\partial f}{\partial \phi} \symbfit{e}_\phi \right) \Delta r \Delta \theta \Delta \phi.
    \end{equation}

    将其代入梯度的积分表示中，并除以体积微元 $\Delta V = r^2 \sin\theta \Delta r \Delta \theta \Delta \phi$，最终推导出球坐标系下的梯度公式：
    \begin{equation}
        \boxed{\operatorname{grad} f = \frac{\partial f}{\partial r} \symbfit{e}_r + \frac{1}{r} \frac{\partial f}{\partial \theta} \symbfit{e}_\theta + \frac{1}{r \sin\theta} \frac{\partial f}{\partial \phi} \symbfit{e}_\phi}
    \end{equation}
\end{solution}

\begin{problem}{高斯公式与调和函数性质}{Gauss-Theorem-Harmonic}
    设函数 $u(x,y,z)$ 在光滑曲面 $S$ 所围成的闭区域 $V$ 上具有直到二阶的连续偏微商，而且满足拉普拉斯方程：
    \begin{equation}
        \Delta u = \frac{\partial^2 u}{\partial x^2} + \frac{\partial^2 u}{\partial y^2} + \frac{\partial^2 u}{\partial z^2} = 0.
    \end{equation}
    试证明：
    \begin{enumerate}
        \item $\oiint_S \frac{\partial u}{\partial \symbfit{n}} \, \d S = 0$；
        \item $\oiint_S u \frac{\partial u}{\partial \symbfit{n}} \, \d S = \iiint_V (\nabla u)^2 \, \d V$，其中 $\frac{\partial u}{\partial \symbfit{n}}$ 是 $u$ 沿 $S$ 外侧法向量 $\symbfit{n}$ 的方向微商。
    \end{enumerate}
\end{problem}

\begin{myproof}
    根据方向导数与梯度的关系，函数 $u$ 沿曲面 $S$ 外侧单位法向量 $\symbfit{n}$ 的方向微商可以表示为：
    \begin{equation} \label{eq:directional-derivative-1}
        \frac{\partial u}{\partial \symbfit{n}} = \nabla u \cdot \symbfit{n}
    \end{equation}

    \ansitem{1}
    证明如下：

    利用方程 \zcref{eq:directional-derivative-1} 以及高斯公式（散度定理），将面积分转化为体积分：
    \begin{equation}
        \begin{aligned}
            \oiint_S \frac{\partial u}{\partial \symbfit{n}} \, \d S & = \oiint_S (\nabla u) \cdot \symbfit{n} \, \d S \\
                                                                     & = \iiint_V \operatorname{div}(\nabla u) \, \d V \\
                                                                     & = \iiint_V \Delta u \, \d V
        \end{aligned}
    \end{equation}
    根据题意，在区域 $V$ 内函数 $u$ 满足拉普拉斯方程 $\Delta u = 0$。代入上式计算，即可得证：
    \begin{equation}
        \boxed{\oiint_S \frac{\partial u}{\partial \symbfit{n}} \, \d S = 0}
    \end{equation}

    \ansitem{2}
    证明如下：

    将被积函数中的方向导数使用方程 \zcref{eq:directional-derivative-1} 展开，并将 $u \nabla u$ 视作一个整体的向量场，应用高斯公式：
    \begin{equation}
        \oiint_S u \frac{\partial u}{\partial \symbfit{n}} \, \d S = \oiint_S (u \nabla u) \cdot \symbfit{n} \, \d S = \iiint_V \operatorname{div}(u \nabla u) \, \d V
    \end{equation}
    接下来计算向量场 $u \nabla u$ 的散度，利用向量微分恒等式展开：
    \begin{equation}
        \begin{aligned}
            \operatorname{div}(u \nabla u) & = \nabla \cdot (u \nabla u)                           \\
                                           & = \nabla u \cdot \nabla u + u (\nabla \cdot \nabla u) \\
                                           & = (\nabla u)^2 + u \Delta u
        \end{aligned}
    \end{equation}
    由已知条件 $\Delta u = 0$ 可知，该散度化简为：
    \begin{equation}
        \operatorname{div}(u \nabla u) = (\nabla u)^2
    \end{equation}
    将其代回前面的体积分表达式中，即完成证明：
    \begin{equation}
        \boxed{\oiint_S u \frac{\partial u}{\partial \symbfit{n}} \, \d S = \iiint_V (\nabla u)^2 \, \d V}
    \end{equation}
\end{myproof}

\endinput
